\documentclass[12pt, letterpaper, parskip = full]{article}
\usepackage[mmddyyyy]{datetime}

\usepackage[utf8]{inputenc}
\usepackage{amsmath}
\usepackage{mathpazo}
\usepackage{eulervm}
\usepackage[skip=10pt plus1pt]{parskip}

\usepackage[margin=1in]{geometry}
%\linespread{1.20}
\usepackage{enumerate}
\usepackage{enumitem}
\usepackage{listings}
\usepackage[T1]{fontenc}
\usepackage{amsfonts}
\usepackage{amssymb}
\usepackage{amsthm}
\usepackage{multirow,array}
\usepackage{graphicx}
\usepackage{physics}
\usepackage{bm}
\usepackage{tikz}
\usepackage{breqn}
\usepackage[longnamesfirst]{natbib}
\usepackage{comment}
\usepackage{setspace}
\usepackage{hyperref}
\usepackage{accents}
\usepackage{graphicx}

\graphicspath{{Images/}}
\setlength{\parindent}{0pt}

\newtheorem{theorem}{Theorem}
\newtheorem{corollary}{Corollary}
\newtheorem{lemma}{Lemma}
\newtheorem{definition}{Definition}
\newtheorem{axiom}{Axiom}
\newcommand{\e}{\mathrm{e}}
\newcommand{\Var}{\mathrm{Var}}
\newcommand{\Ex}[1]{\mathbb{E}\left[#1\right]}
\renewcommand{\Pr}[1]{\mathbb{P}(#1)}
\newcommand{\plim}{\mathrm{plim}}
\newcommand{\pc}[1]{\frac{\dot{#1}(t)}{#1(t)}}
\renewcommand\qedsymbol{$\blacksquare$}
\newcommand{\ubar}[1]{\underaccent{\bar}{#1}}

\title{Homework 5 \\ Actuarial Corporate Finance}
\author{Alejandro Flores-Brown}

\begin{document}

    \maketitle

    \section{Problem 1}

    \subsection{Part 1}

    To find the required return for asset \(B\), we first need to find the risk-free rate. To do this, we will use information about Asset \(A\) and the market portfolio. 

    We are given that the market portfolio has required return \(0.06\). We can guess that under CAPM this should also be the required return since the market portfolio \emph{is} the desired portfolio for every investor, but this can also be shown explicitly with the CAPM equation. This is 
    \[ RR_M = r_f + \beta_M (\Ex{R_M} - r_f) = r_f + \rho_{M,M} \frac{\sigma_M}{\sigma_M} (\Ex{R_M} - r_f) = r_f + \Ex{R_M} - r_f = \Ex{R_M}   \]
    With this in hand, we can now write down the system for the expected and required returns of A:
    \begin{align*}
    \Ex{R_A} &= r_f + \rho_{A,M} \frac{\sigma_A}{\sigma_M} (\Ex{R_M} - r_f) \\
    RR_A &= r_f + \rho_{A,M} \frac{\sigma_A}{\sigma_M} (\Ex{R_A} - r_f)
    \end{align*}
    Substituting values, we can write the first equation as
    \[ \Ex{R_A} = r_f - 0.25 \times \frac{0.5}{0.25} (0.06 - r_f) = r_f - 0.03 + 0.5 r_f = 1.5 r_f - 0.03 \]
    which makes the second equation
    \[ RR_A = 0.03 = r_f - 0.25 \times \frac{0.5}{0.25} (1.5 r_f - 0.03 - r_f) = r_f - 0.25 r_f + 0.015  \]
    Thus we derive
    \[ r_f = \frac{0.03 - 0.015}{0.75} = 0.02 \]
    We then find the expected return of \(R_A\) immediately as \(\Ex{R_A} = 0.02 \times 1.5 - 0.03 = 0\), which we will use in the next part. For now, we look to determine the required return of Asset \(B\).

    The expected return of \(B\) is
    \[ \Ex{B} = r_f + \rho_{B,M} \frac{\sigma_B}{\sigma_M} (\Ex{R_M} - r_f) = 0.02 + 0.3 \times \frac{0.35}{0.25} (0.06 - 0.02) = 0.0368   \]
    which then gives the required return as
    \[  \boxed{RR_B = 0.02 + 0.3 \times \frac{0.35}{0.25} \times (0.368 - 0.02) = 0.027}  \]

    \subsection{Part 2}

    With the above information, it should be easy to derive the expected return of the portfolio. The expected return of the portfolio is just a linear combination of the returns of each asset:
    \[ \Ex{R_P} = 0 \times 0.25 + 0.0368 \times 0.75 = 0.0276   \]
    This then lets us use the CAPM to get
    \[ 0.0276 = 0.02 + \beta{P_M} (0.06 - 0.02) \implies \boxed{\beta_{P,M} = \frac{0.076}{0.04} = 1.9}   \]

    \section{Problem 2}

    The market portfolio is essentially just trying to replicate the return of the market. Since the expected return of the market is just a proportional linear combination of the assets, the market portfolio held by an investor should replicate the weights of the market. That means that we should have
    \[ \frac{shares_{A,M}}{shares_{B,M}} = \frac{shares_{A,P}}{shares_{B,P}} \implies \frac{10000}{5000} = \frac{50}{shares_{B,P}}  \]
    so doing this math, we see that
    \[ \boxed{shares_{B,P} = 25}   \]

\end{document}